\chapter{Preface: A Library is Not a Folder}

Let's be honest with each other. You have a lot of PDFs.

Some of them are books you actually read. Some are papers you downloaded
with good intentions three years ago and have not opened since. Some are
scan artifacts with no metadata, named things like
\code{scan0047-compressed-FINAL.pdf} by a process you no longer remember.
They live in various folders on various machines, and some of them are on
Dropbox, and some are not, and you are not entirely certain which ones are
which.

A folder is not a library. A library has a catalog. It knows what it
contains, who wrote it, what it's about, and where to find it. When you
want something, you don't browse --- you search. When you share a reading
list with a collaborator, you share a manifest, not a zip file.

\textbf{shelf} is a system for making that the reality instead of the aspiration.

It consists of four tools. \code{shelf} is the main one: it owns the
database, manages the files, handles collections and checkouts, exports
bibliographies. \code{shelf-classify} is the metadata brain --- given a
PDF, it figures out what it is, where it fits, and what it should be
called. \code{shelf-index} is the search engine --- it reads your files,
builds a semantic index, and answers natural-language queries against it.
\code{shelf-latex} is the writing integration --- it lets you write
\code{\textbackslash shelf\{Anderson\}} in a \LaTeX{} document and get back
a proper citation with a clickable link that opens the book to the right
page.

The four tools talk to each other using a simple JSON protocol over
standard input and output. \code{shelf} orchestrates; the others are
backends. You can run them together or wire up replacements.

A word about scope. The first version of shelf is aimed squarely at
mathematics: it classifies papers into AMS subject areas, looks up metadata
on zbMATH and arXiv, and proposes filenames in the style mathematicians
use. That specificity is a feature right now --- it's what makes the
defaults useful rather than generic. But the architecture is not
math-specific. Collections, checkouts, semantic search, and clickable
citations work the same way whether your library is algebraic geometry
papers or social science monographs. As helpers for other disciplines are
added, the classification and lookup layers will broaden. The rest of the
manual applies regardless.

This manual is organized into four parts, one per tool. Part~I covers
\code{shelf}: the core concepts, how to manage a library, how collections
and checkouts work, and how to search. Part~II covers \code{shelf-classify}:
how a PDF gets catalogued, what confidence levels mean, and how to call the
classifier directly. Part~III covers \code{shelf-index}: how semantic search
works under the hood and what the three commands do. Part~IV covers
\code{shelf-latex}: the \code{\textbackslash shelf} syntax, the
transformation pipeline, and the fallback package for when you just need
things to compile.

You don't need to read all four parts. If you're setting up a library for
the first time, Part~I is the one. If you're writing a paper and want
clickable citations, jump to Part~IV. If something in the classification
results looks wrong and you want to understand why, Part~II will explain it.
